\section*{Chapter \ref{ch:g9:genes-and-dna} --- Genes and DNA}

\begin{solution}{exo:g9:genes-and-dna:1}
A \omterm{def:g9:genes-and-dna:gene}{gene} is a segment of a \omterm{def:g9:chromosomes:chromosomes}{chromosome} carrying one \omterm{def:g9:heredity-and-traits:traits}{trait}'s
instructions, at a fixed address; \omterm{def:g9:genes-and-dna:gene}{alleles} are its versions
--- the blood-group \omterm{def:g9:genes-and-dna:gene}{gene}'s A, B and O.
\end{solution}

\begin{solution}{exo:g9:genes-and-dna:2}
Because \omterm{def:g9:chromosomes:chromosomes}{chromosomes} come in pairs: one copy of each \omterm{def:g9:genes-and-dna:gene}{gene}
rides each partner --- one from the mother, one from the
father, at the same address.
\end{solution}

\begin{solution}{exo:g9:genes-and-dna:3}
The order of the chemical letters along the strand. Four
kinds of letter, about three billion of them in one human
deck.
\end{solution}

\begin{solution}{exo:g9:genes-and-dna:4}
The strands are letter-by-letter complements: pulled
apart, each rebuilds its partner, and one ladder becomes
two identical ladders --- a full copy per daughter \omterm{def:g5:first-look-at-cells:cell}{cell}.
\end{solution}

\begin{solution}{exo:g9:genes-and-dna:5}
A rare copying error --- changed letters creating a new
\omterm{def:g9:genes-and-dna:gene}{allele}. Mostly neutral, sometimes harmful, occasionally
useful.
\end{solution}

\begin{solution}{exo:g9:genes-and-dna:6}
All \omterm{def:g6:exploring-our-environment:components}{living things} write and read their \omterm{def:g9:genes-and-dna:gene}{genes} in the same
\omterm{def:g9:genes-and-dna:dna}{DNA} letters and code. Consequences: the strongest evidence
for common origin --- and the workability of moving \omterm{def:g9:genes-and-dna:gene}{genes}
between \omterm{def:g5:biodiversity:species}{species}, since the language matches everywhere.
\end{solution}

\begin{solution}{exo:g9:genes-and-dna:7}
Each parent carries brown with blue --- brown expressed,
blue silent. Each deals blue with one chance in two; the
deals are independent, so blue-with-blue lands with one
chance in four --- the blue-eyed child.
\end{solution}

\begin{solution}{exo:g9:genes-and-dna:8}
Group A: A with A, or A with O. Group B: B with B, or B
with O. Group AB: A with B --- both expressed. Group O: O
with O.
\end{solution}

\begin{solution}{exo:g9:genes-and-dna:9}
Both hold the whole library --- the same 46 threads. The
\omterm{def:g8:nervous-communication:neuron}{neuron} reads the nerve-machinery \omterm{def:g9:genes-and-dna:gene}{genes} and \omterm{def:g1:plants-around-us:parts}{leaves} the
\omterm{def:g1:the-five-senses:sense}{skin}'s keratin chapters closed; the \omterm{def:g1:the-five-senses:sense}{skin} \omterm{def:g5:first-look-at-cells:cell}{cell} reads the
reverse. Specialization is a borrowing policy, not a
different library.
\end{solution}

\begin{solution}{exo:g9:genes-and-dna:10}
Mash the \omterm{prop:g3:flowers-fruits-seeds:transform}{fruit} (\omterm{def:g5:first-look-at-cells:cell}{cells} opened); salty water and washing-up
liquid (\omterm{def:g6:cell-unit-of-life:parts}{membranes} --- \omterm{def:g5:first-look-at-cells:cell}{cell}'s and \omterm{def:g6:cell-unit-of-life:parts}{nucleus}'s envelopes ---
\omterm{def:g7:oxygen-in-the-water:dissolved}{dissolved}, \omterm{def:g9:genes-and-dna:dna}{DNA} freed and kept \omterm{def:g7:oxygen-in-the-water:dissolved}{dissolved}); strain (debris
out); cold spirit layered on top (\omterm{def:g9:genes-and-dna:dna}{DNA} undissolves at the
boundary): whitish spoolable threads --- the molecule
itself.
\end{solution}

\begin{solution}{exo:g9:genes-and-dna:11}
Family level: the chin recurs along reproduction's lines,
carried unshown through one generation. \omterm{def:g9:chromosomes:chromosomes}{Chromosome} level:
its determinant rides one thread, halved into \omterm{def:g8:sexual-reproduction:gametes}{gametes} and
restored at \omterm{def:g5:flower-to-fruit:steps}{fertilization}, dealt to some children and not
others. Molecule level: it is an \omterm{def:g9:genes-and-dna:gene}{allele} --- a spelling of
some face-shaping \omterm{def:g9:genes-and-dna:gene}{gene} --- copied faithfully down the
century since the \omterm{prop:g9:genes-and-dna:explains}{mutation} that first wrote it.
\end{solution}

\begin{solution}{exo:g9:genes-and-dna:12}
\omterm{prop:g9:genes-and-dna:explains}{Mutation} writes: copying errors create new spellings ---
molecule level, this chapter. The shuffle deals: halving
and \omterm{def:g5:flower-to-fruit:steps}{fertilization} hand each child a fresh combination ---
\omterm{def:g9:chromosomes:chromosomes}{chromosome} level, \cref{ch:g9:chromosomes} (and
\cref{prop:g8:sexual-reproduction:shuffle}). Expression
shows: of the two copies dealt, what is displayed and what
rides silent --- this chapter's two-copy rule. Only
\omterm{prop:g9:genes-and-dna:explains}{mutation} creates genuinely new spellings; the other two
recombine and reveal what \omterm{prop:g9:genes-and-dna:explains}{mutation} once wrote.
\end{solution}

\begin{solution}{pb:g9:genes-and-dna:1}
\textbf{1.} A \omterm{prop:g9:genes-and-dna:explains}{mutation}: a copying error in the pigment
\omterm{def:g9:genes-and-dna:gene}{gene}, once, in some ancestor's \omterm{def:g8:sexual-reproduction:gametes}{gamete} line --- letters
changed, the instructions spoiled, and the misspelling
copied faithfully ever since.

\textbf{2.} Because P's working instructions build the
pigment regardless of the silent partner: p rides
\emph{carried unshown} --- the album's word, now
molecular.

\textbf{3.} With p on both partners, no working
instructions exist at either address: the pigment
machinery --- the \omterm{def:g9:genes-and-dna:gene}{gene}'s product --- is simply never
built. The \omterm{def:g9:heredity-and-traits:traits}{trait} is the absence of a working spelling.

\textbf{4.} That their pigment \omterm{def:g9:genes-and-dna:gene}{genes} are recognizably the
same \omterm{def:g9:genes-and-dna:gene}{gene} --- similar sequences in the same code ---
kin-copies of one ancestral text, broken independently in
each \omterm{def:g5:biodiversity:species}{species}' own history.

\textbf{5.} The four deals: P-with-P, P-with-p, p-with-P,
p-with-p --- equally likely. The unpigmented build:
one chance in four.

\textbf{6.} The skip: two carriers displaying pigment
transmitted what they did not show, and the hidden
spellings met --- \cref{ex:g9:heredity-and-traits:skip}
run at full mechanism.

\textbf{7.} All carry one p --- their mother's whole
contribution at that address is p. None show the
condition: the partner's P builds the pigment in every
one.

\textbf{8.} Grandmother's p rode one \omterm{def:g9:chromosomes:chromosomes}{chromosome} of a
pair; her \omterm{def:g8:sexual-reproduction:gametes}{gamete}'s halving dealt that partner into the
\omterm{def:g2:animals-grow-up:born}{egg} or sperm that made the carrier parent; the parent's
halving dealt it again into the \omterm{def:g8:sexual-reproduction:gametes}{gamete} that made the
grandchild --- thread to \omterm{def:g8:sexual-reproduction:gametes}{gamete} to \omterm{prop:g8:sexual-reproduction:core}{zygote}, twice over.

\textbf{9.} Family: two ordinary-looking parents, a
one-in-four child, carriage proven by the \omterm{def:g2:born-grow-die:cycle}{birth}.
\omterm{def:g9:chromosomes:chromosomes}{Chromosome}: one pair's partners carrying P and p, halved
and recombined. Molecule: one \omterm{def:g9:genes-and-dna:gene}{gene}, one working and one
misspelled \omterm{def:g9:genes-and-dna:gene}{allele}, the spelling difference deciding
everything.

\textbf{10.} From every judgment: in a carrier, p's
brokenness is invisible --- the partner's P covers it, so
nothing about the carrier's build ``reports'' the
misspelling. What cannot be seen cannot be repaired
away; hidden spellings persist indefinitely.

\textbf{11.} Her \omterm{def:g9:heredity-and-traits:traits}{trait} (no pigment) is inherited ---
p-with-p --- but its \emph{consequence} is lived in an
\omterm{def:g6:exploring-our-environment:components}{environment} full of sun: the missing pigment was the
body's sunscreen, so the \omterm{def:g6:exploring-our-environment:components}{environment} writes on her more
than on her shielded parents --- inheritance setting the
range, the world filling it in.

\textbf{12.} What copies: the double strand, each half
rebuilding its partner at every \omterm{def:g5:first-look-at-cells:cell}{cell} division and
generation. What mis-copied once: a few letters of the
pigment \omterm{def:g9:genes-and-dna:gene}{gene}, in one ancient \omterm{def:g8:sexual-reproduction:gametes}{gamete}. Why the error
persists: the photocopier is faithful --- it copies
whatever is there, misspellings with the same fidelity as
sense.
\end{solution}
